\chapter{Counting Processes and Poisson Calibration}
\label{ch:poisson}

This chapter is the foundation of Part~III. It fixes the mathematical objects --- counting
process, filtration, conditional intensity, compensator --- states the two theorems that
make intensity models \emph{estimable and testable}, and builds the simplest useful model
family: Poisson processes with macro covariates, in continuous and in weekly time. The
through-lines of Chapter~\ref{ch:problem} recur throughout: \emph{identifiability}
(covariate effects are pinned down only by temporal contrast) and \emph{uncertainty}
(over-dispersion is the rule, and every standard error we report is dispersion-corrected).
The chapter ends at the exact boundary of the Poisson world --- clustering beyond
covariates --- which Chapter~\ref{ch:hawkes} crosses.

\section{Counting processes, intensities, compensators}
\label{sec:poi-setup}

A \emph{counting process} $N(t)$, $t \ge 0$, is a right-continuous step function with
$N(0) = 0$ and jumps of size $+1$ at event times $t_1 < t_2 < \cdots$; we assume
simplicity (no coincident points). In our application $N_s(t)$ counts deals in sector $s$
and $N_i(t)$ counts deals of firm $i$. The filtration $\Hist_t$ collects everything
observable strictly before $t$: past events, their marks, and point-in-time covariate
vintages (Chapter~\ref{ch:data}).

\begin{definition}[Conditional intensity]
\label{def:poi-intensity}
The conditional intensity of $N$ with respect to $\Hist_t$ is
\begin{equation}
\lambda(t) \;=\; \lim_{\dd t \downarrow 0}\;
\frac{\Prob\big(N(t + \dd t) - N(t) = 1 \mid \Hist_t\big)}{\dd t},
\label{eq:poi-intensity}
\end{equation}
assumed to exist and to be predictable (determined by information strictly before $t$).
\end{definition}

The intensity is the model: every process in Part~III is specified by writing down
$\lambda(t)$ as a function of time, covariates, and history. Its integral
\begin{equation}
\Lambda(t) \;=\; \int_0^t \lambda(u)\dd u
\label{eq:poi-compensator}
\end{equation}
is the \emph{compensator}: the unique predictable increasing process such that
$M(t) = N(t) - \Lambda(t)$ is a (local) martingale, a property we use without proof
\cite{bremaud1996,daley2003}. Its practical content is unbiasedness of counts against the
model, which is what makes residual analysis on count increments meaningful.

\section{Two foundational theorems}
\label{sec:poi-theorems}

\begin{theorem}[Watanabe characterization]
\label{thm:poi-watanabe}
Let $N$ be a simple point process adapted to $\Hist_t$ with continuous compensator
$\Lambda$. Then $N$ is an inhomogeneous Poisson process --- increments over disjoint
intervals independent, with $N(b) - N(a) \sim \mathrm{Poisson}(\Lambda(b) - \Lambda(a))$
--- if and only if $\Lambda$ is deterministic \cite{daley2003}.
\end{theorem}

Read as a modelling statement: \emph{Poisson means history-free}. The moment $\lambda(t)$
depends on past events, the process is no longer Poisson; conversely, any intensity that
is a deterministic function of time and exogenous covariates inherits the entire Poisson
toolkit of independent increments and closed-form likelihoods.

\begin{theorem}[Time rescaling]
\label{thm:poi-rescaling}
Let $N$ have event times $t_1 < t_2 < \cdots$ and conditional intensity $\lambda(t) > 0$
with $\Lambda(t) \to \infty$ almost surely. Then the transformed times
$\tau_k = \Lambda(t_k)$ are the event times of a unit-rate homogeneous Poisson process;
equivalently, the rescaled inter-arrivals $\delta_k = \tau_k - \tau_{k-1}$ are i.i.d.\
$\mathrm{Exp}(1)$ \cite{ogata1988}.
\end{theorem}

Theorem~\ref{thm:poi-rescaling} holds for \emph{any} point process with an intensity, not
just Poisson --- Hawkes, inhibited, marked, all of them. It is the universal
goodness-of-fit device: substitute the fitted $\hat{\Lambda}$, compute
$\hat{\delta}_k = \hat{\Lambda}(t_k) - \hat{\Lambda}(t_{k-1})$, and test.

\begin{corollary}[Everything in Part~III is testable]
\label{cor:poi-testable}
Every model class in this manual specifies $\lambda(t;\theta)$ explicitly; hence every
fitted model yields rescaled inter-arrivals that should be i.i.d.\ $\mathrm{Exp}(1)$ under
correct specification. Our standing test battery: (i) Kolmogorov--Smirnov test of
$u_k = 1 - e^{-\hat{\delta}_k}$ against $\mathrm{Uniform}(0,1)$; (ii) autocorrelation of
the $u_k$ (independence, which the KS test does not check); (iii) the same by sector and
regime. Failing (ii) while passing (i) means the right intensity level with the wrong
dynamics: the signature of ignored clustering.
\end{corollary}

\section{The homogeneous Poisson process}
\label{sec:poi-homog}

With $\lambda(t) \equiv \lambda$ on an observation window $[0, T]$, the log-likelihood is
$\ell(\lambda) = n \log \lambda - \lambda T$ for $n$ observed events, maximized at
\begin{equation}
\hat{\lambda} \;=\; n / T,
\label{eq:poi-homog-mle}
\end{equation}
with the exact (not asymptotic) confidence interval from the chi-square duality,
\begin{equation}
\lambda \;\in\; \Big[\tfrac{1}{2T}\,\chi^2_{2n,\,\alpha/2},\;\;
\tfrac{1}{2T}\,\chi^2_{2n+2,\,1-\alpha/2}\Big],
\label{eq:poi-homog-ci}
\end{equation}
where $\chi^2_{d,q}$ is the $q$-quantile of the chi-square with $d$ degrees of freedom.
The homogeneous model is acceptable only over windows short enough that macro conditions
and sector heat are effectively constant --- a claim Corollary~\ref{cor:poi-testable}
makes checkable. We use it as a per-window baseline whose rejection is informative.

\section{Inhomogeneous Poisson with covariates}
\label{sec:poi-inhomog}

The first serious model: a deterministic log-linear intensity driven by the macro state,
\begin{equation}
\lambda(t) \;=\; \exp\!\big(a + \beta^{\top} X(t)\big),
\label{eq:poi-loglinear}
\end{equation}
with $X(t)$ piecewise constant at weekly resolution --- the as-of macro vintages of
Chapter~\ref{ch:data}: policy rate, CPI year-on-year, unemployment, each held at its last
published value. Let weeks $w = 1, \dots, W$ partition $[0,T]$ with lengths $\Delta_w$ and
covariate values $X_w$, and let $n_w$ be the event count in week $w$. The log-likelihood
of an inhomogeneous Poisson process,
$\sum_k \log \lambda(t_k) - \int_0^T \lambda(u)\dd u$, becomes
\begin{equation}
\ell(a, \beta) \;=\; \sum_{w=1}^{W} \Big[ n_w \big(a + \beta^{\top} X_w\big)
\;-\; \Delta_w \exp\!\big(a + \beta^{\top} X_w\big) \Big].
\label{eq:poi-loglik}
\end{equation}
Under \eqref{eq:poi-loglinear}, timing within a week carries no information beyond $n_w$:
weekly counts are sufficient, licensing the discrete-time form of
Section~\ref{sec:poi-discrete}.

\begin{proposition}[Strict concavity]
\label{prop:poi-concave}
Write $x_w = (1, X_w^{\top})^{\top}$. The map $(a,\beta) \mapsto \ell(a,\beta)$ in
\eqref{eq:poi-loglik} is concave, and strictly concave if and only if
$\{x_w : w = 1,\dots,W\}$ spans $\mathbb{R}^{p+1}$. The MLE, when it exists, is therefore
unique.
\end{proposition}

\begin{proof}[Proof sketch]
Each term is a linear function of $(a,\beta)$ minus $\Delta_w$ times the exponential of a
linear function; exponentials of linear maps are convex, so $\ell$ is concave. The Hessian
is $-\sum_w \Delta_w e^{a + \beta^{\top} X_w} x_w x_w^{\top}$, negative definite exactly
when no nonzero vector is orthogonal to all $x_w$ --- the spanning condition. This is the
standard exponential-family/log-sum-exp argument.
\end{proof}

Newton iterations converge in a few steps. For inference, with $H$ the observed
information and $\hat{S}$ the outer product of weekly score contributions (clustered by
week, or by sector--week blocks in the multi-sector model), we report the sandwich
\begin{equation}
\widehat{\Var}(\hat{a}, \hat{\beta}) \;=\; H^{-1} \hat{S} H^{-1},
\label{eq:poi-sandwich}
\end{equation}
which remains valid as a quasi-likelihood variance when the Poisson assumption fails in
the ways Section~\ref{sec:poi-overdispersion} documents.

\subsection{Identifiability: covariate effects live in temporal contrast}
\label{sec:poi-ident}

\begin{proposition}[Information equals weighted covariate variance]
\label{prop:poi-ident}
The Fisher information of \eqref{eq:poi-loglik} is
$I(a,\beta) = \sum_w \Delta_w \lambda_w\, x_w x_w^{\top}$ with
$\lambda_w = e^{a + \beta^{\top} X_w}$. Profiling out $a$, the information for $\beta$ is
the intensity-weighted temporal covariance of $X$:
\begin{equation}
I_{\beta} \;=\; \sum_w \omega_w \big(X_w - \bar{X}_{\omega}\big)
\big(X_w - \bar{X}_{\omega}\big)^{\top},
\qquad
\omega_w = \Delta_w \lambda_w,\quad
\bar{X}_{\omega} = \frac{\sum_w \omega_w X_w}{\sum_w \omega_w}.
\label{eq:poi-fisher}
\end{equation}
In particular: (i) a covariate constant over the window is not identified separately from
$a$; (ii) a covariate that is a near-linear trend is near-collinear with any slow trend in
the baseline, and its effective information is the variance of its \emph{detrended}
residual.
\end{proposition}

The proof is the Schur complement of the $(a,a)$ block in $I(a,\beta)$; we omit it. The
consequences matter more. Macro series are slow --- rates move in staircase steps,
unemployment drifts --- so a two-to-three-year window may contain only one or two genuine
covariate contrasts, near-confounded with the baseline level and any secular trend in deal
activity. Our standing guidance: (i) include an explicit trend term so that $\beta$ is
measured against it rather than silently absorbing it; (ii) enter slow series in
differences (rate \emph{changes}, CPI \emph{changes}), not levels; (iii) report effects as
finite contrasts $\beta^{\top}(x_{\mathrm{hi}} - x_{\mathrm{lo}})$ between configurations
that occur in the window, with sandwich intervals; (iv) monitor the condition number of
\eqref{eq:poi-fisher} --- ill-conditioning is an identifiability warning, not a numerical
nuisance.

On the shipped testbed this is not hypothetical: in regime~B the weekly-GLM covariate
estimates recover the daily-world truth with correlation $0.37$ (single seed; protocol in
Chapter~\ref{ch:selection}). Held-out fit still improves substantially over the
covariate-free baseline: predictions and contrasts are far better identified than the
coordinates of $\beta$. Report accordingly.

\section{Discrete-time counts: the weekly GLM}
\label{sec:poi-discrete}

For sector-level work we pass to weekly buckets directly. With $Y_{s,w}$ the count of
sector-$s$ deals in week $w$,
\begin{equation}
Y_{s,w} \;\sim\; \mathrm{Poisson}(\Lambda_{s,w}),
\qquad
\log \Lambda_{s,w} \;=\; a_s + \beta_s^{\top} X_w,
\label{eq:poi-glm}
\end{equation}
independently across $s$ and $w$ given covariates --- a Poisson GLM with log link, one
intercept and slope vector per sector, estimated by the likelihood \eqref{eq:poi-loglik}
summed over sectors. By the sufficiency remark above, when the continuous-time intensity
is constant within weeks this is the same model, not an approximation; feedback enters in
Section~\ref{sec:poi-feedback}. Figure~\ref{fig:poi-sector-counts} shows the raw material.

\begin{figure}[htbp]
\centering
\includegraphics[width=.85\textwidth]{fig_sector_counts}
\caption{Weekly deal counts for three synthetic sectors. Common macro-driven movements are
visible across sectors (the case for shared covariates $X_w$ with sector-specific
loadings $\beta_s$), superimposed on sector-specific levels and bursts that
covariates will turn out not to explain (Section~\ref{sec:poi-overdispersion}).}
\label{fig:poi-sector-counts}
\end{figure}

\section{Over-dispersion}
\label{sec:poi-overdispersion}

The Poisson law imposes $\Var(Y) = \E[Y]$. Deal counts violate it, always, and the
violation is informative.

\begin{definition}[Dispersion]
\label{def:poi-dispersion}
For a fitted count model with means $\hat{\mu}_j$ over $n$ cells and $p$ parameters, the
Pearson dispersion is
\begin{equation}
\hat{\phi} \;=\; \frac{1}{n - p} \sum_{j=1}^{n}
\frac{(y_j - \hat{\mu}_j)^2}{\hat{\mu}_j},
\label{eq:poi-dispersion}
\end{equation}
with $\hat{\phi} \approx 1$ under correct Poisson specification and $\hat{\phi} > 1$
indicating over-dispersion (computed on held-out cells to avoid rewarding overfit means).
\end{definition}

Two distinct mechanisms inflate $\hat{\phi}$, and both are present in deal data.
\emph{Unobserved heterogeneity}: if $Y \mid Z \sim \mathrm{Poisson}(Z\mu)$ with a latent
positive factor $Z$, $\E[Z] = 1$ (risk appetite, unmodelled sector health), then
$\Var(Y) = \mu + \mu^2 \Var(Z) > \mu$. Taking $Z$ gamma-distributed yields the negative
binomial as a mixture Poisson --- we state this without derivation and use it when a full
likelihood is needed. \emph{Contagion}: if events raise the intensity of subsequent events
(Section~\ref{sec:poi-feedback}, Chapter~\ref{ch:hawkes}), counts cluster and the
marginal variance again exceeds the mean. The cheap, robust correction is quasi-Poisson:
keep the point estimates from \eqref{eq:poi-glm} and multiply all standard errors by
$\hat{\phi}^{1/2}$ (consistent with the sandwich \eqref{eq:poi-sandwich}).

\begin{figure}[htbp]
\centering
\includegraphics[width=.85\textwidth]{fig_dispersion}
\caption{Weekly variance versus mean per sector, against the Poisson benchmark
$\Var = \E$ (dashed). Every sector sits above the line. On regime-B held-out weeks the
fitted covariate GLM \eqref{eq:poi-glm} leaves a Pearson dispersion of $2.72$: conditional
on the macro covariates, weekly counts are nearly three times more variable than Poisson.
All GLM standard errors in this manual are quasi-Poisson corrected accordingly.}
\label{fig:poi-dispersion}
\end{figure}

\begin{warning}[Dispersion does not identify its cause]
\label{warn:poi-dispersion}
A held-out dispersion of $2.72$ tells us the covariate-conditional Poisson model is wrong;
it does not tell us \emph{how}. Latent common factors and genuine contagion produce the
same first-order symptom --- variance above the mean and positively autocorrelated
residual counts. Quasi-Poisson repairs the standard errors under either mechanism, which
is why it is our default; but any \emph{scientific} claim ("deals beget deals") requires
separating the two, and that separation cannot be done from dispersion statistics. It
needs timestamp-level structure, and it is the central identification question of
Chapter~\ref{ch:hawkes}.
\end{warning}

\section{Excitation and inhibition in count GLMs}
\label{sec:poi-feedback}

The GLM extends naturally to history dependence: append lagged counts to the predictor,
\begin{equation}
\log \Lambda_{s,w} \;=\; a_s + \beta_s^{\top} X_w
\;+\; \sum_{\ell=1}^{L} b_{\ell}^{\top} Y_{w-\ell},
\label{eq:poi-feedback}
\end{equation}
where $Y_{w-\ell}$ stacks all sectors' counts, so $b_{\ell}$ captures within- and
cross-sector feedback at lag $\ell$. By Theorem~\ref{thm:poi-watanabe} this is no longer a
Poisson process --- the compensator is now random --- but conditional on the past each
$Y_{s,w}$ is still Poisson, the likelihood factorizes, and estimation is unchanged.

Two features make \eqref{eq:poi-feedback} the right first model for history effects.
First, \emph{negative coefficients are legal}: the log link keeps $\Lambda_{s,w} > 0$ for
any sign of $b_{\ell}$, so \emph{inhibition} --- deal flow suppressing near-future deal
flow through capital absorption or cooldowns --- is expressed as naturally as excitation.
The linear Hawkes process of Chapter~\ref{ch:hawkes} cannot do this: its intensity is a
sum of nonnegative kernels, and negative excitation would violate positivity. Second, the
weekly clock makes stability auditable --- with a subtlety the following analysis makes
precise.

\paragraph{Stability: two mechanisms.}
Collect the feedback matrices $B_{\ell}$ ($M \times M$, sectors) and let
$G = \sum_{\ell} B_{\ell}$.

\emph{Mechanism A (additive feedback).} If the feedback enters the intensity linearly,
$\Lambda_w = c + \sum_{\ell} B_{\ell} Y_{w-\ell}$ with nonnegative entries, the process is
a discrete-time branching (INAR-type) system: each event contributes on average $G$
offspring events, and $\rho(G) < 1$ --- spectral radius below one --- suffices for
stationarity, with stationary mean $(I - G)^{-1} c$. This is the direct analogue of the
Hawkes branching-ratio condition of Chapter~\ref{ch:hawkes}.

\emph{Mechanism B (log link on raw counts).} Model \eqref{eq:poi-feedback} feeds raw
counts through an exponential: the conditional-mean map is
$\mu(y) = \exp(a + B y)$ elementwise. This is \emph{geometric}, not additive, feedback. A
useful local summary is the Jacobian at the deterministic fixed point $\mu^{*}$,
\begin{equation}
G_{\mathrm{eff}} \;=\; \mathrm{diag}(\mu^{*})\, G,
\qquad
\rho_{\mathrm{eff}} \;=\; \rho(G_{\mathrm{eff}}),
\label{eq:poi-rhoeff}
\end{equation}
the \emph{effective spectral radius}. But $\rho_{\mathrm{eff}} < 1$ guarantees only local
stability of the mean map: the Poisson conditional law has unbounded support, and a
sufficiently large fluctuation can push counts past the basin of attraction, after which
the exponential feedback compounds and the process diverges. Global ergodicity for
log-linear count autoregressions requires genuinely stronger conditions, developed by
\cite{fokianos2009}; the standard safe designs are to feed $\log(1 + Y_{w-\ell})$ instead
of raw counts (making the feedback additive on the log scale, Mechanism~A-like), to cap
counts, or to keep $\rho_{\mathrm{eff}}$ well below one.

\begin{warning}[A measured metastability episode]
\label{warn:poi-metastable}
This failure mode is not theoretical. The regime-A synthetic generator originally used
log-link feedback on raw counts with effective radius $\rho_{\mathrm{eff}} = 0.35$ ---
comfortably subcritical by the local criterion. The configuration was \emph{metastable}:
across most weeks it behaved normally, then a single burst crossed the basin boundary and
tipped a world calibrated to roughly $50\mathrm{k}$ events into $169\mathrm{k}$ events.
Rebuilt at $\rho_{\mathrm{eff}} = 0.15$, the generator is stable (Chapter~\ref{ch:selection}
documents the episode). The lesson generalizes: \emph{never} simulate --- or deploy
forecasts from --- log-link raw-count feedback near criticality; keep excitation terms
small and let covariates carry the signal, or move the feedback to the $\log(1+Y)$ scale
where Mechanism~A logic applies.
\end{warning}

\begin{remark}[Fitted feedback overstates contagion under aggregation]
\label{rem:poi-aggregation}
The estimated $\hat{b}_{\ell}$ in \eqref{eq:poi-feedback} must not be read as a measured
contagion strength. On the testbed, a daily-world generator with \emph{modest} true
contagion and a $7$-day excitation decay yields a fitted weekly-GLM effective spectral
radius of $0.57$. Two mechanisms conspire: weekly binning aliases within-week and
adjacent-week excitation into the lag-one coefficient (an aggregation artifact), and
latent-quality clustering --- persistent hot sectors --- loads onto lagged counts because
they are the best available proxy for the latent factor. Treat fitted feedback as an upper
bound on contagion and as a forecasting device; the identification of true excitation is
deferred to Chapter~\ref{ch:hawkes}.
\end{remark}

\section{Simulation}
\label{sec:poi-simulation}

For piecewise-constant deterministic intensities, simulation is trivial: draw independent
Poisson counts per week (and uniform timestamps within weeks if needed). History-dependent
intensities require thinning \cite{lewis1979}.

\paragraph{Algorithm (thinning).}
Given a predictable intensity $\lambda(t)$ and, at each time, a computable local bound
$\bar{\lambda} \ge \lambda(u)$ valid on a lookahead interval:
\begin{enumerate}
\item Set $t \leftarrow 0$; initialize history $\Hist$; compute the current bound
$\bar{\lambda}$ valid on $[t, t + \Delta)$.
\item Draw $E \sim \mathrm{Exp}(\bar{\lambda})$ and propose $t' = t + E$. If
$t' > \min(T, t + \Delta)$, advance $t$ to that endpoint, refresh $\bar{\lambda}$, and
repeat from step~2; stop when $t \ge T$.
\item Draw $U \sim \mathrm{Uniform}(0,1)$; accept an event at $t'$ if
$U \le \lambda(t') / \bar{\lambda}$, evaluating $\lambda(t')$ with the history up to
$t'$.
\item If accepted, append $t'$ to $\Hist$ (excitation raises future $\lambda$, so refresh
$\bar{\lambda}$). Set $t \leftarrow t'$ and return to step~2.
\end{enumerate}
Correctness requires only that $\bar{\lambda}$ dominate $\lambda$ on its validity
interval. Note the interaction with Section~\ref{sec:poi-feedback}: for log-link feedback
on raw counts no global bound exists and local bounds explode near bursts --- one more
reason that family is dangerous to simulate.

\paragraph{Discrete-time simulation with risk-set caps.}
Our production simulators run on the weekly clock: for $w = 1, \dots, W$ compute
$\Lambda_{s,w}$ from covariates and lagged counts, draw
$Y_{s,w} \sim \mathrm{Poisson}(\Lambda_{s,w})$, then cap at the risk set,
$Y_{s,w} \leftarrow \min(Y_{s,w}, |\RS_{s,w}|)$, where $\RS_{s,w}$ is the set of firms
currently eligible to transact (alive, not in a post-deal cooldown,
Chapter~\ref{ch:hawkes}). The cap is substantive --- a finite firm universe cannot
generate unbounded weekly counts --- and doubles as an explosion guard for feedback
models. We log every cap activation: a simulator that hits its caps more than rarely is
reporting a misspecified intensity, and the fix belongs in the model, not in a larger cap.

\section{What Poisson cannot capture}
\label{sec:poi-limits}

The covariate GLM \eqref{eq:poi-glm} says: given the macro path, disjoint weeks are
independent and variance equals the mean. The data answer with held-out dispersion $2.72$
and positively autocorrelated rescaled residuals --- clustering of events beyond anything
the covariates explain. Within this chapter's model class the only response is defensive:
quasi-Poisson inflation of uncertainty, feedback terms read as forecasting devices.

The scientific question underneath is sharper. Two generative stories produce exactly
these symptoms: a \emph{latent common factor} --- unobserved risk appetite multiplying
every sector's intensity, inducing correlated bursts with no causal link between deals ---
and \emph{self-excitation} in the sense of \cite{hawkes1971}, where each deal genuinely
raises the intensity of successors via signaling, competitive urgency, and follow-on
capital. The two stories have the same second moments at
the weekly layer (Remark~\ref{rem:poi-aggregation} is this equivalence measured), yet they
imply different counterfactuals: under contagion, seeding a deal raises subsequent
activity; under a latent factor, it changes nothing. Distinguishing them requires
timestamp-level structure, explicit kernels, and honest statements about what remains
unidentified. That is Chapter~\ref{ch:hawkes}.
